% TOP_PROBLEM: {"id": 2768, "title": "Kirby Problem 2.20", "category_id": 7, "category": {"id": 7, "name": "topology", "display_name": "Topology", "description": "Properties preserved under continuous deformations.", "slug": "topology", "order_index": 7, "created_at": "2026-07-31T15:26:25.671Z"}, "status": "open"}
% FIRST_POSTED: null
% ENTRY_KIND: statement_only
% Imported from: batch11.tex; UnsolvedMath ID 2768
\documentclass[11pt]{article}
\usepackage[margin=25mm]{geometry}
\usepackage{fontspec}
\setmainfont{Latin Modern Roman}
\setsansfont{Latin Modern Sans}
\usepackage{amsmath,amssymb,mathtools}
\usepackage{unicode-math}
\setmathfont{Latin Modern Math}
\usepackage{xurl,hyperref}
\setcounter{secnumdepth}{0}
\setlength{\parindent}{0pt}
\setlength{\parskip}{5pt}
\setlength{\emergencystretch}{3em}
\newcommand{\sref}[2]{\hyperlink{source:#1:#2}{[S#2]}}
\newcommand{\cataloguescope}[1]{\paragraph{Recorded target.}#1}
\begin{document}
% UnsolvedMath ID 2768; source catalogue code KP-2.20
\setcounter{section}{2767}
\section{Kirby Problem 2.20}\label{2768}
{\small\sffamily UnsolvedMath ID 2768 · Topology}
\subsection{Definitions and mathematical statement}

\paragraph{Shared notation.}$\mathbb N=\{1,2,\ldots\}$, and $\mathbb Z,\mathbb Q,\mathbb R,\mathbb C$ denote the integers, rationals, reals and complex numbers. Any other conventions are specified locally.


Let $S_g^b$ be a compact oriented genus-$g$ surface with $b$ boundary components, and $\operatorname{Mod}(S_g^b)$ its orientation-preserving mapping class group fixing boundary pointwise. Write $t_c$ for the right-handed Dehn twist about a simple closed curve $c$ and $[a,b]=aba^{-1}b^{-1}$. Monodromy of a genus-$g$ Lefschetz fibration over a genus-$h$ surface is encoded by a relation
\[t_{c_1}\cdots t_{c_\ell}[a_1,b_1]\cdots[a_h,b_h]=1\]
in $\operatorname{Mod}(S_g)$, with positive twists about essential vanishing cycles at the critical values. Boundary-marked lifts can have the product of boundary twists on the right; for a pencil these boundary components encode its base points and the base genus is $h=0$. A surface bundle is the case $\ell=0$. A Lefschetz fibration has isolated critical points locally modeled on $(z_1,z_2)\mapsto z_1z_2$; a pencil is defined away from finitely many base points with local ratio model $z_1/z_2$. Holomorphicity here requires integrable complex structures on the total space and base, making the fibration or bundle holomorphic and the pencil meromorphic with these base points. An almost-complex structure alone does not suffice.
\paragraph{Statement recorded in the supplied source.}
Determine properties of the monodromy factorization that detect whether the represented Lefschetz pencil, Lefschetz fibration, or surface bundle over a surface admits such a holomorphic realization. In particular, characterize the extra restrictions satisfied by monodromy factorizations of holomorphic fibrations, for all three classes of maps. \sref{2768}{2}

\subsection{Short English statement}
Find conditions on monodromy factorizations that distinguish integrable holomorphic pencils, fibrations and surface bundles from their more general topological or symplectic counterparts.
\subsection{Sources}
\begin{description}
\item[\hypertarget{source:2768:1}{S1}] ULAM.AI and the UnsolvedMath Contributors, UnsolvedMath: A Curated Collection of Open Mathematics Problems, record 2768 (KP-2.20). CC BY 4.0. \url{https://huggingface.co/datasets/ulamai/UnsolvedMath}.
\item[\hypertarget{source:2768:2}{S2}] R. İnanç Baykur, Robion C. Kirby and Daniel Ruberman, eds., \emph{K3: A New Problem List in Low-Dimensional Topology}, Mathematical Surveys and Monographs 295, American Mathematical Society (2026), Problem 2.20 and its remarks; Chapters 1 and 2 for the stated conventions. \url{https://math.berkeley.edu/sites/default/files/surv-295-ruberman-watermarked-author-pdf.pdf}.
\end{description}
\label{end:2768}
\end{document}
